\documentclass[a4paper]{article}
\usepackage{titlesec}
\usepackage{geometry}
\geometry{hmargin={2cm,2cm}, vmargin={3.5cm,3.5cm}}
\usepackage{graphicx}
\usepackage[table]{xcolor}
\usepackage{hyperref}
\hypersetup{pdfborder = {0 0 0},
colorlinks,
citecolor=black,
filecolor=black,
linkcolor=black,
urlcolor=blue}
\definecolor{logo}{RGB}{20,99,163}
\newcommand{\sectionbreak}{\clearpage}
\usepackage[T1]{fontenc}
\usepackage[utf8]{inputenc}
\usepackage[english,italian]{babel}
\usepackage{fancyhdr}
\usepackage{lastpage}

% creazione comando subsubsubsection
\titleclass{\subsubsubsection}{straight}[\subsection]

\newcounter{subsubsubsection}[subsubsection]
\renewcommand\thesubsubsubsection{\thesubsubsection.\arabic{subsubsubsection}}
\titleformat{\subsubsubsection}
  {\normalfont\normalsize\bfseries}{\thesubsubsubsection}{1em}{}
    \titlespacing*{\subsubsubsection}
    {0pt}{3.25ex plus 1ex minus .2ex}{1.5ex plus .2ex}

\makeatletter
\def\toclevel@subsubsubsection{4}
\def\l@subsubsubsection{\@dottedtocline{4}{7em}{4em}}
\makeatother

\setcounter{secnumdepth}{4}
\setcounter{tocdepth}{4}
%fine creazione comando


%Comandi di impaginazione uguale per tutti i documenti
\pagestyle{fancy}
\lhead{\includegraphics[scale=0.2]{../../../Images/logo.png}}
%Titolo del documento
\rhead{\large Verbale esterno 2021/03/17}
%\rfoot{\thepage}
\cfoot{Pagina \thepage{} di \pageref{LastPage}}
\setlength{\headheight}{35pt}
\renewcommand{\footrulewidth}{0.4pt}




\begin{document}
\begin{titlepage}
    \begin{center}
        \includegraphics{../../../Images/logo.png}\\
        \vspace{20px}
        \textcolor{logo}{\hrulefill}\\
        \vspace{20px}
        \textbf{\huge\textcolor{logo}{Verbale esterno 2021/03/17}}\\
        \vspace{10px}
        \textcolor{logo}{\hrulefill}\\
        \vspace{40px}
        \textbf{\Large Informazioni sul documento}\\
        \vspace{20px}
        \begin{tabular}{p{100px} | p{100px}}
            \textbf{Versione} & 0.0.1\\
            \textbf{Responsabile} & Responsabile\\
            \textbf{Redattori} & Giaccone Nicolò \newline *Redattore2*\\
            \textbf{Verificatore} & Verificatore\\
            \textbf{Destinatari} & Prof. Tullio Vardanega \newline Prof. Riccardo Cardin, \newline TechSWEave\\
            \textbf{Stato} & Da approvare\\
            \textbf{Uso} & Esterno\\
        \end{tabular}\\
        \vspace{60px}
        \href{mailto:techsweave@gmail.com}{techsweave@gmail.com}\\

    \end{center}
    \end{titlepage}

    %DIARIO DELLE MODIFICHE
    \begin{center}
        \textbf{\Large Diario delle modifiche}\\
        \vspace{10px}
        \begin{table}[h!]
        \centering
        \rowcolors{2}{logo!10}{logo!40}
        \renewcommand{\arraystretch}{1.8}
        \begin{tabular}{p{150px} p{90px} p{50px} p{60px} p{45px}}
            \rowcolor{logo!70} \textbf{Modifica} & \textbf{Autore} & \textbf{Ruolo} & \textbf{Data} & \textbf{Versione}\\
            Prima stesura del verbale & Nicolò Giaccone & Analista & 2021-03-18 & 0.0.1 \\    
        \end{tabular}
    \end{table}
    \end{center}

    %INDICE
    \newpage
    \tableofcontents
    \newpage
    % CORPO DEL DOCUMENTO
    \section{Informazioni generali}
    \begin{itemize}
        \item Data e ora: 2021/03/17 ore 17:30-18:15;
        \item Luogo: Google meet fornito dai proponenti;
        \item Partecipanti interni:  
        \begin{itemize}
            \item  Samuele De Simone;
            \item  Manuel Varo;
            \item  Tito Scutari;
            \item  Nicolò Giaccone;
            \item  Amedeo Meggiolaro;
            \item  Simone Urbani;
            \item  Marco Barbaresco;
        \end{itemize}
        \item Partecipanti esterni: 
        \begin{itemize}
            \item Alessandro Maccagnan; 
            \item Milo Ertola;
        \end{itemize} 
        \item Segretario: Nicolò Giaccone.
    \end{itemize}
\section{Argomenti della riunione}
\begin{itemize}
    \item Presentazione del gruppo all'azienda;
    \item Discussione sulle tempistiche della realizzazione del progetto didattico;
    \item Chiarimenti su alcuni aspetti del progetto;
    \item Proposta di alcuni requisiti;
    \item Discussione sull'approccio da seguire per svolgere al meglio il progetto;
    \item Discussione sulle criticità della consegna;
    \item Chiarimenti sulla formulazione dei documenti.
\end{itemize}
\section{Resoconto}
\subsection{Presentazione del gruppo all'azienda}
I primi minuti della chiamata sono stati dedicati alla presentazione dei componenti del gruppo TechSWEave ai proponenti.
\subsection{Discussione sulle tempistiche della realizzazione del progetto didattico}
Si è discusso sulle tempistiche di realizzazione del progetto sia in relazione alle consegene fissate per il progetto didattico, si per la mole di lavoro che deve essere sostenuta per la realizzazione del progetto stesso.
\subsection{Chiarimenti su alcuni aspetti del progetto}
Sono state proposte alcuni dubbi sul capitolato, che erano sono sorti durante la stesura dei vari documenti:
\begin{itemize}
    \item La piattaforma consiste di un unico venditore;
    \item Scelta della lingua da utilizzare nella piattaforma;
\end{itemize}

\end{document}
